\section{Introduction}

When faced with uncertainty, biological neural systems don't deliberate
indefinitely---they act on the first confident signal. This principle, refined
by millions of years of evolution, underlies rapid decision-making from prey
detection to social judgment. In cortical winner-take-all circuits, competing
neural populations race to threshold; the first to fire suppresses alternatives
and triggers action~\citep{grossberg1987competitive}. The speed of convergence
itself carries information: a fast response indicates a clear, unambiguous
stimulus.

We apply this biological insight to improve neural network ensembles. Standard
ensemble methods average predictions across models, treating all outputs
equally regardless of when each model reaches its answer. For iterative
reasoners with adaptive computation time (ACT)~\citep{graves2016act}, this
discards valuable information: models that converge quickly have typically
found ``clean'' solution paths, while those still deliberating are often stuck
or uncertain.

Our key observation is that \textbf{inference speed is an implicit confidence
signal}. When multiple models reason in parallel, the first to halt is most
likely correct. This simple selection rule---which we call \emph{halt-first
ensembling}---improves accuracy by 5.7 percentage points over probability
averaging while using 10$\times$ less compute on the Sudoku-Extreme benchmark.

This raises a natural question: can we internalize this benefit within a single
model? Training and deploying multiple models is expensive. We show that the
diversity exploited by halt-first selection comes primarily from different
random initializations during training. By maintaining $K$ parallel latent
initializations within one model and training with a winner-take-all (WTA)
objective, we achieve ensemble-level accuracy with single-model deployment
cost.

\paragraph{Contributions.}
\begin{enumerate}
  \item We demonstrate that halt-first selection outperforms probability
    averaging for ensembles of iterative reasoners, achieving 97.2\% accuracy
    vs.\ 91.5\% with 10$\times$ less compute.
  \item We introduce \emph{oracle-first training}, a WTA method that
    internalizes ensemble diversity within a single model, achieving
    96.9\% $\pm$ 0.6\% accuracy with zero inference overhead.
  \item We identify chain convergence during training and introduce
    \emph{expunging}---removing converged hypotheses to recover compute for
    new samples.
  \item We discover that Muon optimization fails with standard SwiGLU and
    provide a fix: layer normalization before the down-projection. This
    enables 4$\times$ higher learning rates than AdamW.
  \item We provide a theoretical interpretation connecting halting speed to
    the spectral properties of bilevel equilibria, grounding ``speed is
    confidence'' in the conditioning of nested fixed points.
  \item We document extensive hyperparameter exploration across 180+
    experiments, identifying critical factors: carry policy, SVD-aligned
    initialization, and z\_L diversity. Code and trained models at
    \url{https://anonymous.4open.science/r/TODO}.
\end{enumerate}
